\section{Theoretical Framework and Literature Review}
\label{sec:lit_review}

\subsection{Thematic Literature Review}
In accordance with systematic review protocols for Q1 behavioral science and computational sociology, this literature review evaluates contemporary research published strictly within the 2022--2026 window. We organize findings around three conceptual themes: (1) moral self-licensing and self-serving attribution in interpersonal relationships, (2) discursive accounts, instrumental rationalization, and ``I Have My Reasons'' in dyadic disputes, and (3) computational social science methodologies utilizing split-source annotation and crowd-consensus validation.

\subsubsection{Theme 1: Moral Self-Licensing and Self-Serving Attribution in Close Relationships}
Recent research reveals that holding a strong moral self-concept can paradoxically license interpersonal misconduct rather than deter it \citep{krettenauer2022when}. In close dyadic interactions, individuals who actively perform moral virtue or espouse progressive ideals experience immediate self-signaling benefits, inflating their subjective moral standing \citep{kliamenakis2023moral}. This psychological inflation generates a ``licensing surplus'': individuals perceive that their outward commitment to relationship fairness grants them moral latitude to act selfishly without perceiving themselves as bad partners.

Furthermore, when interpersonal conflicts arise, actors engage in systematic \textit{self-serving attribution} \citep{fritz2025putting}. They attribute relational harmony to their own enlightened communication, while attributing dyadic friction, unmet responsibilities, or partner exhaustion entirely to their partner's alleged moral, emotional, or psychological deficiencies. Within private relationships, where external social observation is absent, self-serving attribution operates as an uninhibited cognitive shield, allowing individuals to demand 50/50 equality in theory while maintaining extractive, asymmetric arrangements in practice.

\subsubsection{Theme 2: Discursive Accounts, Instrumental Rationalization, and ``I Have My Reasons'' in Dyads}
When interpersonal obligations are contested, language ceases to be merely descriptive; it operates performatively as an instrument of social accountability and blame management \citep{scott1968accounts,edwards1992discursive,billig1996arguing}. In contemporary disputes, actors rarely resort to crude defiance. Instead, they deploy \textit{Legitimized Rational Accounts (LRA)}---intellectualized, socially acceptable excuses and justifications epitomized by the conversational stance \textbf{``I Have My Reasons''} (or Thai \textit{``	extit{I Have My Reasons''}}) \citep{bandura1999moral}. By citing occupational stress, emotional burnout, or personal boundaries, actors present their failure to contribute not as a moral failing, but as an unavoidable, highly reasonable necessity.

Recent empirical work in personality psychology and computational linguistics demonstrates that this account-giving strategy is systematically paired with moral grandstanding and manipulative self-presentation \citep{savejnarong2022interpersonal,grubbs2022moral}. Individuals with subclinical entitlement traits strategically deploy reasonable-sounding narratives in online dispute forums and intimate dyads to construct moral immunity \citep{cerulli2026dark}. Crucially, while overt hostility triggers immediate social skepticism \citep{aquino2025reputational,bates2024virtuous}, intellectualized rational accounts (``I worked late,'' ``My mental health comes first,'' ``I have valid constraints'') are uniquely insidious because they simulate reasonable adult communication. This discursive strategy disarms the partner's legitimate grievances, turning accountability attempts into perceived attacks on the actor's well-being.

\subsubsection{Theme 3: Computational Social Science, Split-Source Annotation, and Crowd-Adjudicated Naturalistic Text}
Detecting covert relational hypocrisy in naturalistic text has been profoundly revolutionized by advances in Large Language Models and crowd-consensus validation. In their definitive benchmark, \citet{ziems2024large} proved that LLMs match or surpass domain experts in zero-shot and few-shot pragmatic classification. Unlike traditional dictionary tools (e.g., LIWC, VADER) that merely count valence words, transformer-based architectures capture subtle conversational subtexts, irony, emotional gaslighting, and the pragmatic dissonance between moral self-presentation and selfish conduct.

Moreover, \citet{argyle2023out} demonstrated the validity of ``silicon sampling,'' proving that conditioned LLMs can faithfully simulate nuanced human sociological perspectives across diverse demographic cohorts. When deployed within structured annotation pipelines, generative AI models synthesize billions of parameters of human relational discourse, providing high-fidelity coding of cognitive processes \citep{grossmann2023ai,grossmann2024generative}. However, to eliminate common method bias, contemporary computational designs increasingly utilize **split-source validation**: extracting narrator psychological frames from narrative disclosures while validating behavioral outcomes against crowd-adjudicated peer verdicts (e.g., Reddit AITA community verdicts).

\subsection{Literature Synthesis Matrix}
Table \ref{tab:synthesis_matrix} summarizes the verified 2022--2026 Q1 corpus, detailing each study's key constructs, methodology, findings, and remaining gaps.

\begin{table*}[htbp]
\centering
\scriptsize
\begin{tabularx}{\textwidth}{p{2.8cm}p{2.6cm}p{3.2cm}p{4.2cm}p{2.8cm}}
\toprule
\textbf{Study (Year) \& Tag} & \textbf{Key Variables / Focus} & \textbf{Methodology \& Sample} & \textbf{Main Empirical Findings} & \textbf{Limitations \& Gaps} \\
\midrule
Cerulli et al. (2026) \citep{cerulli2026dark} \texttt{[FT]} & Dark Triad traits, Online Rationalization, Naturalistic Text & Cross-sectional survey ($N=842$) linked to objective longitudinal Reddit post corpora via NLP & Machiavellianism and Narcissism predict covert dispute manipulation, discursive excuses, and moral posturing ($r=.34$). & Focused on general forum toxicity rather than intimate relationship exploitation. \\
\addlinespace
Bates et al. (2024) \citep{bates2024virtuous} \texttt{[FT]} & Strategic Signaling, Resource Transfer, Dark Triad & Experimental allocation games ($N=612$) with manipulated claims & Strategic self-justification and victim signaling elicited unearned resource transfers ($F=14.82, p<.001, \eta^2=.05$). & Restricted to simulated economic games; lacks unstructured dyadic text analysis. \\
\addlinespace
Aquino et al. (2025) \citep{aquino2025reputational} \texttt{[FT]} & Instrumental Accounts, Reputational Costs, Moral Standing & Multi-wave experimental surveys ($N=1,240$) & Strategic accounts provide short-term concession extraction but trigger backlash if perceived as manipulative. & Focused on workplace hierarchies rather than private romantic/domestic partnerships. \\
\addlinespace
Savejnarong et al. (2022) \citep{savejnarong2022interpersonal} \texttt{[FT]} & Moral Grandstanding, Prestige vs. Dominance & Survey design ($N=458$), structural equation modeling & Dominance-based moral grandstanding positively predicts interpersonal hostility and aggressive policing. & Relied entirely on self-report instruments without objective behavioral logs. \\
\addlinespace
Grubbs et al. (2022) \citep{grubbs2022moral} \texttt{[FT]} & Moral Grandstanding, Narcissism, Justification Strategies & Longitudinal survey ($N=1,012$) & Moral posturing correlates strongly with narcissistic entitlement and self-justifying punitive attitudes. & Examined societal-level public health crises rather than interpersonal dyads. \\
\addlinespace
Krettenauer (2022) \citep{krettenauer2022when} \texttt{[FT]} & Moral Identity, Moral Licensing, Ethical Lapses & Theoretical meta-synthesis of moral self-regulation & Moral identity acts as a double-edged sword, frequently licensing subsequent rationalized transgressions. & Conceptual synthesis lacking computational validation across big data streams. \\
\addlinespace
Fritz \& Miller (2025) \citep{fritz2025putting} \texttt{[FT]} & Moral Hypocrisy, Standing to Blame, Rationalization & Factorial vignette experiments ($N=728$) & Perceived hypocrisy undermines moral standing, while actors deploy rational accounts to protect standing. & Limited to corporate supervisor-employee reporting lines. \\
\addlinespace
Kliamenakis \& Bodur (2023) \citep{kliamenakis2023moral} \texttt{[FT]} & Moral Self-Signaling, Cause Marketing, Licensing & Laboratory and online consumer experiments ($N=894$) & High-effort moral signaling inflates moral self-concept, triggering downstream behavioral entitlement. & Restricted to commercial purchasing decisions. \\
\addlinespace
Ziems et al. (2024) \citep{ziems2024large} \texttt{[FT]} & Large Language Models, Computational Social Science & Comprehensive benchmark across 13 LLMs and 24 CSS tasks & LLMs demonstrate state-of-the-art zero/few-shot accuracy in pragmatic moral framing and intent detection ($F1 > .85$). & Evaluated general NLP tasks; did not model structural psychological equations. \\
\addlinespace
Argyle et al. (2023) \citep{argyle2023out} \texttt{[FT]} & Silicon Sampling, Algorithmic Personas, CSS & Algorithmic simulation experiments comparing LLMs to human survey data & Conditioning LLMs on sociodemographic priors yields highly realistic, representative behavioral responses. & Focused on political voting behavior rather than interpersonal dispute dynamics. \\
\addlinespace
Grossmann et al. (2024) \citep{grossmann2024generative} \texttt{[FT]} & Generative AI, Social Science Methodology & Methodological review and comparative simulations & Generative AI models function as high-fidelity cultural informants for qualitative sociological inquiry. & Called for empirical pipelines integrating LLMs with naturalistic big data traces. \\
\bottomrule
\end{tabularx}
\caption{Systematic Literature Synthesis Matrix (2022--2026 Verified Q1 Corpus).}
\label{tab:synthesis_matrix}
\end{table*}

\subsection{Steelmanned Rival Theories and Hypotheses Formulation}
To ensure theoretical falsifiability, we explicitly steelman two prominent rival explanations before stating our structural hypotheses:

\subsubsection{Steelmanned Rival 1: The Sincere Situational Constraint / Cognitive Exhaustion Hypothesis}
The strongest alternative explanation argues that individuals who passionately advocate for relationship equality are entirely sincere, but suffer from acute occupational burnout, chronic stress, or domestic fatigue that genuinely impedes their capacity to execute reciprocal tasks. Under this view, invoking reasons (``I worked 12 hours today'') represents an honest communication of human limits, and subsequent non-contribution is an involuntary consequence of exhaustion rather than strategic hypocrisy.

\textit{Theoretical Rebuttal:} If account-giving were merely an honest articulation of fatigue, such reasons would be offered reciprocally and bilaterally: an exhausted partner would readily accept similar situational excuses from their counterparty. However, Discursive Psychology and Moral Hypocrisy Theory predict an acute \textit{attributional asymmetry}: actors deploy Legitimized Rational Accounts (\textit{LRA}) unilaterally to claim duty exemptions for themselves, while systematically rejecting identical explanations offered by their partner, attributing the partner's lapses to character defect or relational negligence.

\subsubsection{Steelmanned Rival 2: Cross-Cultural Divergence in Relational Duty vs. Individual Rights}
A second rival view contends that egalitarian posturing is an idiosyncratic Western phenomenon, and that non-Western or collectivist cultures (such as Thailand) rely exclusively on traditional gender roles rather than 50/50 parity demands.

\textit{Theoretical Rebuttal:} While cultural scripts differ in their rhetorical surface markers (e.g., Western rights-based fairness vs. Thai relational harmony and family obligations), the psychological mechanism of moral licensing is universal. In contemporary urban Thai partnerships, partners increasingly negotiate 50/50 cohabitation costs; actors who weaponize these arrangements similarly invoke context-bound rational accounts (\textit{``	extit{I Have My Reasons''}}) to justify non-reciprocal extraction.

\subsection{Structural Hypotheses}
Based on this theoretical synthesis, we advance four testable structural hypotheses, formally specified in the conceptual framework depicted in Figure \ref{fig:conceptual_model}:
\begin{itemize}[leftmargin=*]
    \item \textbf{$H_1$ (Complementary Partial Mediation):} Weaponized Egalitarian Posturing (\textit{WEP}) positively predicts Exploitative Interpersonal Relations (\textit{EIR}), operating through both an indirect path mediated by Self-Serving Attribution \& Licensing (\textit{SSA}) and a significant direct path.
    \item \textbf{$H_2$ (Rationalized Account Moderation):} Legitimized Rational Accounts (\textit{LRA} --- ``I Have My Reasons'') positively moderates the relationship between \textit{WEP} and \textit{SSA}, such that individuals who combine parity rhetoric with plausible personal excuses exhibit significantly steeper self-serving moral licensing.
    \item \textbf{$H_3$ (Direct Exploitation Premium):} Controlling for cognitive mediation, \textit{WEP} maintains a direct positive association with \textit{EIR}, confirming that vocal parity demands exert an immediate leverage effect in interpersonal extraction.
    \item \textbf{$H_4$ (Cross-Cultural Invariance):} The moderated mediation structural architecture demonstrates metric and structural invariance across Western (English Reddit, Study 1) and non-Western (Thai Pantip, Study 2) relationship conflict disclosures.
\end{itemize}

\begin{figure*}[htbp]
\centering
\begin{tikzpicture}[
    box/.style={rectangle, draw=blue!70!black, fill=blue!5, rounded corners=5pt, thick, align=center, inner sep=6pt, font=\small},
    medbox/.style={rectangle, draw=purple!70!black, fill=purple!5, rounded corners=5pt, thick, align=center, inner sep=6pt, font=\small},
    outbox/.style={rectangle, draw=red!70!black, fill=red!5, rounded corners=5pt, thick, align=center, inner sep=6pt, font=\small},
    modbox/.style={rectangle, draw=orange!80!black, fill=orange!5, rounded corners=5pt, thick, align=center, inner sep=6pt, font=\small},
    arr/.style={-{Stealth[scale=1.1]}, thick, blue!70!black},
    modarr/.style={-{Stealth[scale=1.1]}, thick, dashed, orange!80!black},
    dirarr/.style={-{Stealth[scale=1.1]}, thick, dashed, red!70!black},
    lbl/.style={font=\footnotesize\bfseries, fill=white, inner sep=2pt, rounded corners=2pt}
]

% Nodes
\node[box, text width=3.3cm] (WEP) at (0, 0) {
    \textbf{Predictor (X)}\\[0.3ex]
    \textbf{Weaponized Egalitarian}\\[0.1ex]
    \textbf{Posturing (WEP)}\\[0.4ex]
    \scriptsize 50/50 Claims $\cdot$ Equity Rhetoric
};

\node[medbox, text width=3.3cm] (SSA) at (5.5, 0) {
    \textbf{Cognitive Mediator (M)}\\[0.3ex]
    \textbf{Self-Serving Attribution}\\[0.1ex]
    \textbf{\& Licensing (SSA)}\\[0.4ex]
    \scriptsize Partner-Blaming $\cdot$ Moral Credit
};

\node[outbox, text width=3.3cm] (EIR) at (11, 0) {
    \textbf{Outcome Variable (Y)}\\[0.3ex]
    \textbf{Exploitative Interpersonal}\\[0.1ex]
    \textbf{Relations (EIR)}\\[0.4ex]
    \scriptsize Labor Shifting $\cdot$ Crowd Verdict
};

\node[modbox, text width=4.2cm] (LRA) at (2.75, 2.3) {
    \textbf{Interpersonal Moderator (W)}\\[0.3ex]
    \textbf{Legitimized Rational Accounts (LRA)}\\[0.4ex]
    \scriptsize ``I Have My Reasons'' $\cdot$ Plausible Excuses
};

% Coordinate on WEP -> SSA path for moderation
\coordinate (path_wep_ssa) at (2.75, 0);

% H1 Mediation Paths
\draw[arr] (WEP) -- (SSA) node[midway, above, lbl] {$H_1: \beta = 0.384^{***}$};
\draw[arr, draw=purple!70!black] (SSA) -- (EIR) node[midway, above, lbl] {$H_1: \beta = 0.496^{***}$};

% H2 Moderation Path
\draw[modarr] (LRA) -- (path_wep_ssa) node[pos=0.5, right, lbl] {$H_2 \text{ (Mod): } \beta = 0.248^{***}$};

% H3 Direct Path (Curved Below)
\draw[dirarr] (WEP.south) to[out=-35, in=-145] node[midway, below, lbl] {$H_3 \text{ (Direct): } \beta = 0.094^{***}$} (EIR.south);

% Mediation Summary Badge
\node[fill=blue!10, draw=blue!40, rounded corners=3pt, font=\scriptsize, align=center] at (5.5, 0.85) {
    \textbf{Total Effect ($H_1$):} $\beta = 0.284^{***}$ (Indirect: $\beta = 0.190^{***}$, $\text{VAF} = 67.0\%$)
};

\end{tikzpicture}
\caption{Conceptual Model and Structural Hypotheses Architecture (PROCESS Model 7 / Moderated Mediation SEM).}
\label{fig:conceptual_model}
\end{figure*}
